\subsubsection{アーキテクチャにおける「よさ」}
アーキテクチャの「よさ」は、一つの尺度だけでは測れない。ある利害関係者にとってよいアーキテクチャが、別の利害関係者にとっては問題を含むことがある。たとえば、セキュリティを非常に強くした構成は安全に見えるが、費用が高く、検証が難しく、利用者の操作が複雑になることがある。

建築の古典的な考え方では、建物には強さ、用、見た目のよさが必要だとされる。この考え方をソフトウェアに置き換えると、次の問いになる。
\begin{itemize}
  \item 長期運用と進化に耐えられるか。
  \item 意図した用途に合っているか。
  \item このアーキテクチャで、費用対効果よく構築できるか。
  \item 構築、利用、保守を行う人にとって明確で理解しやすいか。
\end{itemize}

\keyterm{アーキテクチャトレードオフ分析法}\en{Architecture Tradeoff Analysis Method, ATAM}は、品質特性に基づいてアーキテクチャを評価する方法である。品質特性を整理し、シナリオを使って評価し、品質要求同士のトレードオフを分析する。ほかにも、ソフトウェアアーキテクチャ分析法\en{SAAM}や品質特性ワークショップ\en{QAW}などの方法がある。

\begin{statementbox}{評価の本質}
評価の目的は、アーキテクチャに点数を付けることだけではない。重大なリスク、隠れた前提、品質特性間の衝突、説明できない判断を早く発見することである。
\end{statementbox}
